\capitulo{1}{Introducción}

Descripción del contenido del trabajo y del estructura de la memoria y del resto de materiales entregados.

Se ha decidido realizar este trabajo por los beneficios que presentaría para el cliente, en este caso, el Campus de la Politécnica de la Universidad de Burgos, el disponer de una aplicación web en la que poder gestionar los horarios del campus, las reservas, el calendario de exámenes de convocatorias y la ocupación de las aulas, todo en una misma aplicación.

Esta aplicación de la que ahora mismo la universidad no dispone de algo parecido, es muy necesario ya que facilita y automatiza mucho, tanto el trabajo de la realización de los horarios de cada campus, como la gestión de reservas, todo ello vinculado a la ocupación de aulas. Y es que no es nada práctico, por ejemplo, tener toda esta información almacenada en Excels, ya que cuando tú mueves por ejemplo, una reserva o una asignatura de día u hora, tienes que ir a la ocupación del aula y cambiarlo tú también manualmente, por el contrario, esta aplicación lo hace de manera automática, disminuyendo el trabajo y por consecuencia minimizar el error humano. Además de que cuando alguien realiza una solicitud de una reserva o directamente se crea una, el sistema nos muestra que aulas se encuentran disponibles para la fecha y hora marcadas, evitando así tener que ir mirando aula por aula la disponibilidad para el plazo especificado.

Además, esta aplicación también nos permite la exportación de los horarios de los distintos grados, los exámenes de convocatoria por grado, la ocupación de las aulas semanal por las asignaturas (reservas periódicas) y las reservas puntuales por aula.